\chapter{Théorie des nombres}


\section{Division euclidienne}

\begin{thm}[Division euclidienne]
	Soient $a, b \in \mathbb{Z}$, $b \neq 0$. Il existe un unique couple $(q, r) \in \mathbb{Z} \times \mathbb{N}$ tel que
	\[
	a = bq + r, \qquad 0 \leq r \leq |b| - 1.
	\]
	Les entiers $q$ et $r$ s'appellent \textit{le quotient} et \textit{le reste de la division euclidienne} de $a$ par $b.$
\end{thm}

\begin{itemize}
	\item \textbf{Divisibilité :} $b$ (non nul) \textit{divise} $a$, noté $b \mid a$, s'il existe un entier $q$ tel que $a = bq$, c'est-à-dire si le reste de la division de $a$ par $b$ est $0.$ On dit aussi que $a$ est un \textit{multiple} de $b$.
\end{itemize}

% ----------------------------------------------------------


\begin{exercise}
	\begin{itemize}
	\item[$\bullet$]Trouver tous les entiers $n \geq 1$ tels que $\dfrac{n^3 + 3}{n^2 + 7}$ soit un entier.



	\item[$\bullet$] Soient $a$ et $b$ deux entiers et $n \geq 1$ un entier. Montrer que $a - b$ divise $a^n - b^n$. Montrer que, si $n$ est impair, $a + b$ divise $a^n + b^n$ (si $a=b$ ou $a=-b$, le diviseur est nul et l'on convient que $0$ divise seulement $0$).
\end{itemize}
\end{exercise}
\subsection*{Nombres premiers}

\begin{defi}
	Soit $p \geq 2$ un entier. Le nombre entier $p$ est \textit{premier} si ses seuls diviseurs positifs sont 1 et $p$.\\
	\quad Un entier $n \geq 2$ est dit \textit{composé} (non premier) si $n = ab$ avec $a \geq 2$ et $b \geq 2$.
\end{defi}

\begin{fact}
	Tout entier $n \geq 2$ admet un diviseur premier.
\end{fact}

\begin{fact}
	Soient $a$ et $b$ deux entiers et $p$ un nombre premier. Si $p \mid ab$, alors $p \mid a$ ou $p \mid b$.
\end{fact}

\begin{fact}[Théorème fondamental de l'arithmétique]
	Tout entier supérieur ou égal à $2$ s'écrit comme un produit de nombres premiers, et cette écriture est unique à l'ordre des facteurs près.
\end{fact}

\begin{exercise}
	\begin{itemize}
		\item[$\bullet$] Montrer qu'un nombre premier supérieur à $3$ est de la forme $6n \pm 1$.
		\item[$\bullet$]Soient $a$ et $b$ deux entiers supérieurs ou égaux à $2$. Montrer que $a^4 + 4b^4$ n'est pas un nombre premier.
	\end{itemize}
\end{exercise}

% ----------------------------------------------------------
\section{PGCD et PPCM}

Soient $a$ et $b$ deux entiers non tous deux nuls. Les diviseurs positifs communs à $a$ et $b$ forment un ensemble non vide (il contient $1$) et fini (si $a \neq 0$, chacun est au plus égal à $|a|$ ; sinon, au plus égal à $|b|$) : il a donc un plus grand élément. Si $a$ et $b$ sont non nuls, les multiples strictement positifs communs à $a$ et $b$ forment une partie non vide de $\mathbb{N}$ (elle contient $|ab|$) : elle a donc un plus petit élément. Ces deux résultats viennent du chapitre « Principe extrémal ». Les définitions suivantes ont donc un sens :
\begin{defi}
	Soient $a$ et $b$ deux entiers tels que $(a, b) \neq (0, 0)$. Le plus grand des diviseurs positifs communs à $a$ et $b$ existe ; on le note $\pgcd(a,b).$ Si $a$ et $b$ sont non nuls, le plus petit des multiples strictement positifs communs à $a$ et $b$ existe ; on le note $\ppcm(a,b).$ Si $\pgcd(a,b)=1,$ on dit que $a$ et $b$ sont \textit{premiers entre eux}.
\end{defi}

\begin{fact}
	Si $(a,b)\neq(0,0)$ : $\pgcd(a, b) = \pgcd(b,\, a - b)$. Si de plus $b \neq 0$ et $a = bq + r$ (division euclidienne), alors $\pgcd(a, b) = \pgcd(b,\, r)$.
\end{fact}

\begin{fact}
	Pour $a, b$ non nuls : $|ab| = \pgcd(a, b) \cdot \ppcm(a, b)$.
\end{fact}

\begin{fact}[Théorème de Bézout]
	Soient $a$ et $b$ deux entiers non tous nuls. Alors il existe deux entiers $x$ et $y$ tels que $\pgcd(a, b) = ax + by$.\\
	\quad Conséquence : $a$ et $b$ sont premiers entre eux si et seulement s'il existe deux entiers $x$ et $y$ tels que $ax + by = 1$. (Sens direct : c'est le théorème. Réciproque : tout diviseur positif commun à $a$ et $b$ divise $ax + by = 1$, donc vaut $1$.)
\end{fact}

\begin{exercise}
	\begin{itemize}
		\item[$\bullet$] Soit $n$ un entier. Calculer $\pgcd(5n + 2,\; 12n + 5)$ et $\pgcd(9n + 4,\; 2n - 1)$ (la réponse peut dépendre de $n$).
	\item[$\bullet$]Soient $a$ et $b$ deux entiers non tous nuls. Soient $c, x, y$ trois entiers tels que $c = ax + by.$ Montrer que $\pgcd(a, b)$ divise $c.$
\end{itemize}
\end{exercise}

% ----------------------------------------------------------
\section{Congruences}

\begin{defi}
	Soit $n \in \mathbb{N}^*$ et soient $a$ et $b$ deux entiers. On dit que $a$ et $b$ sont congrus modulo $n$ (noté $a \equiv b \pmod{n}$) si $a$ et $b$ admettent le même reste après division euclidienne par $n$. Cela revient à dire que $n$ divise $a - b$.
\end{defi}

\begin{fact} Soient $a, b, c$ des entiers et $n \geq 1$. Si $\pgcd(c, n) = 1$ et $ca \equiv cb \pmod{n}$, alors $a \equiv b \pmod{n}$.
\end{fact}

\begin{fact}
	Soient $n \geq 1$ et $f$ un polynôme à coefficients dans $\mathbb{Z}$. Si $a \equiv b \pmod{n}$, alors $f(a) \equiv f(b) \pmod{n}$.
\end{fact}

\begin{fact}[Petit théorème de Fermat]
	Soient $a$ un entier et $p$ un nombre premier. Alors
	\[
	a^p \equiv a \pmod{p}.
	\]
\end{fact}

\begin{exercise}
	\begin{itemize}
		\item[$\bullet$] Montrer que l'équation $15x^2 - 7y^2 = 9$ n'admet pas de solutions entières.
		\item[$\bullet$] Montrer qu'il n'existe pas de polynôme $f$ à coefficients entiers tel que $f(7) = 11$ et
		$f(11) = 13$.
	\end{itemize}
\end{exercise}

% ----------------------------------------------------------
\section{Partie entière}

Soit $x$ un réel. L'ensemble des entiers inférieurs ou égaux à $x$ est non vide et majoré par $x$ ; comme toute partie non vide et majorée de $\mathbb{Z}$, il a un plus grand élément (principe extrémal). La définition suivante a donc un sens :
\begin{defi}
	Soit $x \in \mathbb{R}$. Le plus grand entier inférieur ou égal à $x$ est appelé la \textit{partie entière} de $x$ ; on le note $\lfloor x \rfloor$ ou $E(x).$ Ainsi $\lfloor x \rfloor \leq x < \lfloor x \rfloor + 1.$
\end{defi}

\begin{exercise} Soient $x, y$ deux nombres réels et $n$ un entier naturel non nul.
	\begin{itemize}
		\item[$\bullet$] Montrer : $E(x) + E(y) \leq E(x + y)$.
		\item[$\bullet$] Montrer : $\lfloor \lfloor x \rfloor / n \rfloor = \lfloor x/n \rfloor$.
		\item[$\bullet$] Montrer : $\lfloor x \rfloor + \lfloor y \rfloor + \lfloor x + y \rfloor \leq \lfloor 2x \rfloor + \lfloor 2y \rfloor$.
		\item[$\bullet$] Calculer $\left\lfloor \sqrt{1} \right\rfloor + \left\lfloor \sqrt{2} \right\rfloor + \cdots + \left\lfloor \sqrt{2026} \right\rfloor$.
	\end{itemize}
\end{exercise}
\section{Problèmes}

\begin{problem}[8]
	Soit $A=3^{105}+4^{105}.$ Montrer que $7$ divise $A.$ Calculer le reste de $A$ modulo $11$ et le reste de $A$ modulo $13.$
	
	\end{problem}

\begin{problem}[17]
	Montrer que si $p$ est un nombre premier supérieur ou égal à $5$, on a $p^2\equiv 1\pmod {24}.$
	
\end{problem}

\begin{problem}[97]
	Trouver les trois derniers chiffres de l'écriture décimale de $7^{9999}.$ 
	
\end{problem}

\begin{problem}[118]
	Montrer que l'équation $y^2=x(x+1)(x+2)(x+3)$ n'admet pas de solution $(x,y)$ en entiers strictement positifs. 

\end{problem}







